% 多源BFS：3x5网格（4邻接），障碍格(1,2)不属于图，起点集合 {(0,0),(2,4)}。
% 距离（行0..2）：[0,1,2,3,2] / [1,2,#,2,1] / [2,3,2,1,0]，由 scratchpad 脚本核对。
\begin{tikzpicture}[font=\figurefont\fontsize{8}{10}\selectfont,
  vtx/.style={circle,draw=FigureStroke,line width=.8pt,minimum size=6.5mm,inner sep=0pt,text=FigureInk},
  note/.style={font=\FigureNoteFont\fontsize{7.5}{9}\selectfont,text=FigureMuted}]
  \def\dx{1.05}\def\dy{1.0}
  % 距离表：行 i，列 j
  \foreach \i/\row in {0/{0,1,2,3,2},1/{1,2,-1,2,1},2/{2,3,2,1,0}}{
    \foreach \val [count=\j from 0] in \row {
      \ifnum\val<0\else
        \ifnum\val=0
          \node[vtx,fill=FigureYellowFill] (n\i\j) at (\j*\dx,-\i*\dy) {$\val$};
        \else
          \node[vtx,fill=FigureBlueFill] (n\i\j) at (\j*\dx,-\i*\dy) {$\val$};
        \fi
      \fi
    }
  }
  % 障碍
  \draw[FigureRed,line width=.8pt] (2*\dx-0.25,-\dy-0.25) rectangle (2*\dx+0.25,-\dy+0.25);
  \draw[FigureRed,line width=.8pt] (2*\dx-0.25,-\dy-0.25)--(2*\dx+0.25,-\dy+0.25) (2*\dx-0.25,-\dy+0.25)--(2*\dx+0.25,-\dy-0.25);
  % 边
  \foreach \i in {0,1,2}{\foreach \j/\k in {0/1,1/2,2/3,3/4}{
    \pgfmathtruncatemacro{\bad}{(\i==1 && (\j==2 || \k==2)) ? 1 : 0}
    \ifnum\bad=0 \draw[FigureStroke,line width=.8pt] (n\i\j)--(n\i\k);\fi}}
  \foreach \j in {0,1,2,3,4}{\foreach \i/\k in {0/1,1/2}{
    \pgfmathtruncatemacro{\bad}{(\j==2) ? 1 : 0}
    \ifnum\bad=0 \draw[FigureStroke,line width=.8pt] (n\i\j)--(n\k\j);\fi}}
  \node[anchor=west] at (5.0,-0.2) {黄色：起点集合，距离 0};
  \node[anchor=west] at (5.0,-0.75) {数字：到最近起点的最少边数};
  \node[anchor=west] at (5.0,-1.3) {红框：不属于图的格子};
  \node[anchor=west,note] at (5.0,-1.85) {只从左上角出发时，右下角距离为 6};
\end{tikzpicture}
